\subsection{Consistency Across Visual Features}
While mode grades summarize the most common performance level a simplification technique achieves for preserving visual features across different simplification degrees, they do not indicate how consistently a simplification technique behaves across different datasets. To capture this dimension, we compute the \emph{deviation}---the percentage of datasets that receive grades \emph{different} from the mode---for every simplification technique and visual feature pair.

Formally, if a simplification technique achieves mode grade $M$ on $m$ out of $n$ datasets for a given visual feature, the deviation is:

$$\text{Deviation} = \frac{n - m}{n} \times 100\%$$

This deviation reflects the \emph{data dependency} of each method: whether its performance remains stable across domains or changes dramatically depending on dataset characteristics. Unlike variance-based measures, this percentage-based deviation has a direct, intuitive interpretation: \emph{what fraction of datasets deviate from the typical (mode) grade}.

We classify deviation into three categories based on the proportion of datasets that receive the mode grade:

\begin{itemize}
    \item \textbf{Consistent (Deviation $< 25\%$):} More than 75\% of datasets receive the mode grade, indicating that the simplification technique behaves similarly regardless of data characteristics.
    \item \textbf{Moderately data-dependent (Deviation $25\%$--$50\%$):} 50--75\% of datasets receive the mode grade, suggesting partial dependency on dataset structure.
    \item \textbf{Highly data-dependent (Deviation $> 50\%$):} Less than 50\% of datasets receive the mode grade---excellent on some datasets and poor on others---revealing strong interaction between the simplification technique and data characteristics.
\end{itemize}

\textbf{Interpreting tied mode grades.} When two grades occur with equal frequency, both are considered mode grades. In such cases, deviation is calculated using the frequency of \emph{one} of the tied modes (not their sum). For example, consider a simplification technique where across 80 datasets: 27 receive grade A, 27 receive grade B, and 26 receive other grades (C, D, or F). Here, both A and B are mode grades (tied at 27 datasets each). The deviation is calculated as $\frac{80 - 27}{80} = 66.25\%$, meaning 66.25\% of datasets did not receive the mode grade of 27. This high deviation indicates substantial variability: while A and B each appear in one-third of datasets, the remaining third exhibits entirely different behavior. The tied modes do not indicate consensus but rather polarization in performance.

The complete deviation matrix is reported in \Cref{tab:variance}. Across all 19 simplification techniques and 23 feature preservation metrics, the observed deviations span a wide range, from near-perfect consistency (minimum deviation = 0\%) to extreme variability (maximum deviation = 72.5\%). Of the 437 simplification technique and visual feature pairs, 29.3\% show high consistency (deviation $< 25\%$), 39.8\% fall in the moderately data-dependent range (deviation 25\%--50\%), and 26.8\% exhibit high data-dependency (deviation $> 50\%$).

When examining average deviation across all 23 metrics, simplification techniques cluster into three distinct groups. However, the true insight emerges not from these broad averages but from the \emph{feature-selective} nature of consistency: many methods exhibit dramatically different stability patterns across visual feature categories.

\begin{itemize}
    \item \textbf{Predominantly data-independent methods with selective vulnerabilities.} Six simplification techniques achieve low deviation (below 25\%) for at least half of the 21 feature metrics, indicating that their behavior is largely predictable across diverse datasets. However, each exhibits specific feature categories where data dependency emerges. \methodGau{} demonstrates the strongest overall stability (16 low, 5 medium, 0 high), yet even Gaussian smoothing shows moderate-to-high deviation for extrema and spikes/dips preservation. Five additional methods demonstrate substantial consistency: \methodMean{} (13 low, 7 medium, 1 high), \methodSG{} (11 low, 8 medium, 2 high), \methodDP{} (11 low, 6 medium, 4 high), \methodButter{} (10 low, 8 medium, 3 high), and \methodMed{} (10 low, 5 medium, 6 high). 
    
    All six methods share a critical commonality: moderate-to-high deviation for extrema and spikes/dips preservation, often accompanied by elevated deviation for curvature or periodicity. This systematic pattern reveals that \emph{shape-based features are universally challenging for simplification algorithms}, exhibiting strong data-dependency even when the same algorithms achieve near-perfect stability for level, noise, and trend metrics. Notably, \methodMed{} exhibits the most severe volatility with 6 metrics exceeding 50\% deviation, while \methodMean{} has only 1 such metric despite both achieving similar low-deviation counts.
    
    \item \textbf{Moderately data-dependent methods.} Eight simplification techniques exhibit moderate deviation (25\%--50\%) as their dominant or substantial pattern: IIR filters \methodCheb{} (5 low, 15 medium, 1 high) and \methodEllip{} (4 low, 14 medium, 3 high), frequency-domain filter \methodFFT{} (4 low, 11 medium, 6 high), aggregator \methodPAA{} (1 low, 14 medium, 6 high), downsamplers \methodM4{} (2 low, 14 medium, 5 high) and \methodLTTB{} (9 low, 6 medium, 6 high), hybrid downsampler \methodMinMaxLTTB{} (8 low, 8 medium, 5 high), and subsampler \methodUniform{} (5 low, 10 medium, 6 high). These methods occupy a middle ground: too variable for reliable deployment across arbitrary datasets, yet not unpredictable enough to be unusable. \methodPAA{} now achieves one low-deviation metric (a shift from zero under the previous calculation), yet remains among the most data-dependent aggregators. \methodLTTB{} exhibits balanced deviation across categories (equal medium and high counts), reflecting its visual area-based selection strategy that performs consistently on distributional features but variably on local geometric structures depending on data density. \methodMinMaxLTTB{} shows equal low and medium deviation counts, demonstrating how its hybrid approach of combining LTTB stability with extrema preservation creates moderate variability: reliable for trend and noise metrics but more data-dependent for shape-based features like curvature and spikes/dips.
    
    \item \textbf{Universally data-dependent methods.} Five methods exhibit substantial high deviation: \methodMin{} (0 low, 8 medium, 13 high), \methodFPCS{} (2 low, 2 medium, 17 high), \methodASAP{} (0 low, 5 medium, 16 high), \methodTDA{} (1 low, 8 medium, 12 high), and \methodMax{} (7 low, 7 medium, 7 high). Critically, \methodMin{} and \methodASAP{} have \emph{zero} low-deviation metrics across all 21 features, indicating complete absence of stable preservation guarantees regardless of feature type. \methodFPCS{} remains the most extreme case, exhibiting high deviation for 17 out of 21 metrics (81\%), meaning fewer than one-fifth of its feature preservation outcomes are predictable across datasets. \methodASAP{} follows closely with 16 high-deviation metrics (76\%). This stands in stark contrast to \methodGau{} which has zero high-deviation metrics. \methodMax{} exhibits perfect balance across deviation categories (7-7-7 distribution), reflecting extreme-value filtering's fundamental data-dependency: stable for derivative features but highly variable for level and trend metrics depending on whether data trends upward or contains noise. Similarly, \methodMin{}'s instability reflects the same characteristic: effective on monotonic signals, destructive on noisy data. \methodASAP{}'s poor consistency is particularly noteworthy given its design goal of adaptive optimization---its autocorrelation-based heuristic creates radically different outcomes depending on signal periodicity, making it excellent for periodic data but unreliable for irregular time series.
\end{itemize}

It is important to note that mode grades and deviation describe complementary aspects of simplification technique behavior:
\begin{itemize}
    \item \textbf{Mode grades measure typical quality:} the most common performance level a technique achieves across datasets for a given feature.
    \item \textbf{Deviation measures stability:} what percentage of datasets deviate from that typical performance.
\end{itemize}

A simplification technique with a high mode grade but high deviation (e.g., \methodMed{} on certain visual features) may be excellent for specific dataset characteristics but unreliable elsewhere. Conversely, a technique with low mode grade and low deviation (e.g., \methodDP{}) is consistently poor and therefore not suitable for most tasks. Together, mode grades and deviation provide a comprehensive characterization of the suitability of simplification techniques, distinguishing methods that are reliably good, reliably bad, or highly data-dependent.
